\section{Свойства слабых топологий.}
\begin{definition}
    Пусть $(X, \tau)$ -- некоторое линейное пространство, снабжённое топологией $\tau$.
    Будем называть $\tau$ -- векторной топологией, если $\tau$ удовлетворяет первой аксиоме отделимости и $+\colon X \times X \to X, \ (x, y) \mapsto x + y$ и $\cdot\colon \Cm \times X \to X, \ (\alpha, x) \mapsto \alpha x$ являются непрерывными.
    Их непрерывность эквивалентна непрерывности их в каждой точке:
    \[
        \begin{cases}
            \forall x, y \in X \ \forall U(x + y) \in \tau \ \exists W(x), W(y) \in \tau\colon W(x) + W(y) \subset U(x + y), \\
            \forall \alpha \in \Cm\ \forall x \in X \ \forall U(\alpha x) \in \tau \ \exists B_\epsilon(\alpha) \ \exists W(x) \in \tau \colon B_\epsilon(\alpha)W(x) \subset U(\alpha x).
        \end{cases}
    \]
    Само $(X, \tau)$ будем называть топологическим векторным пространством.
\end{definition}
\begin{note}
    Для топологических векторных пространств первая и вторая аксиомы отделимости эквивалентны.
\end{note}
\begin{note}
    Заметим, что если $(X, \tau)$ -- топологическое векторное пространство, и $G \in \tau$, то $\forall x \in X$ множество $G + x$ также открыто.

    Это очевидно, поскольку трансляция $T_{x}(y) = y + x$ -- непрерывна в силу непрерывности сложения и имеет непрерывную обратную трансляцию $T_{-x}$.
    А значит она переводит открытое в открытое -- в частности, $T_x (G) = G + x$ будет открытым.

    Но это означает то, что топология $X$ определяется лишь системой окрестностей нуля.
\end{note}
\begin{definition}
    Пусть $(X, \tau)$ -- топологическое векторное пространство.
    Локальной базой $\beta_0$ называется такое семейство открытых множеств, что
    \[
        \forall U(0) \in \tau \ \exists W_0 \in \beta_0\colon W_0 \subset U(0).
    \]
\end{definition}
\begin{note}
    Из непрерывности умножения на скаляр следует и то, что умножение на ненулевой скаляр сохраняет открытость, то есть если
    \[
        M_{\lambda}\colon (X, \tau) \to (X, \tau), \ x \mapsto \lambda x,
    \]
    при $\lambda \neq 0$, и $G \in \tau$, то $\lambda G$ также лежит в $\tau$.
    Аналогично, $M_{\frac{1}{\lambda}}$ -- также непрерывно, и является обратным к $M_\lambda$ -- а значит $M_\lambda$ переводит открытые множества в открытые.
\end{note}
\begin{proposition}
    Пусть $X$ -- линейное пространство и $\Phi$ -- семейство функционалов, разделяющих точки на $X$.
    Тогда $\tau_\Phi$ -- векторная топология на $X$.
\end{proposition}
\begin{proof}
    Известно, что $\tau_\Phi$ -- хаусдорфова топология.

    Проверим непрерывность сложения и умножения на скаляр.
    Пусть $x, y \in X$ и $U(x + y) \in \tau_{\Phi}$.
    По определению базы $\exists \phi_1, \ldots \phi_N \in \Phi$ и $\epsilon > 0$:
    \[
        \bigcap_{k = 1}^{N} V(x + y, \phi_k, \epsilon) \subset U(x + y).
    \]
    Рассмотрим окрестности
    \[
        W_1(x) = \bigcap_{k = 1}^{N} V\left(x, \phi_k, \frac{\epsilon}{2}\right), \quad W_2(y) = \bigcap_{k = 1}^{N} V\left(y, \phi_k, \frac{\epsilon}{2}\right).
    \]
    Тогда для любого $\xi + \eta \in W_1(x) + W_2(y)$, где $\xi \in W_1(x)$ и $\eta \in W_2(y)$, то
    \[
        \forall k \ |\phi_k(\xi + \eta) - \phi_k(x + y)| \leq |\phi_k(\xi) - \phi_k(x)| + |\phi_k(\eta) - \phi_k(y)| < \epsilon,
    \]
    и $\xi + \eta \in U(x + y)$.

    Покажем теперь непрерывность умножения на скаляр.
    Пусть $\lambda \in \Cm$, $x \in X$ и $U(\lambda x) \in \tau_\Phi$.
    Тогда
    \[
        \exists \bigcap_{k = 1}^{N} V(\lambda x, \phi_k, \epsilon) \subset U(\lambda x)
    \]
     Теперь надо подобрать такое $\delta > 0$, что для $|\alpha - \lambda| < \delta$ множество
    \[
        \alpha \cdot \bigcap_{k = 1}^{N} V(x, \phi_k, \epsilon) \subset U(\lambda x).
    \]
    Если $y \in \bigcap_{k = 1}^{N} V(x, \phi_k, \delta)$, то
    \[
        |\phi_k(\alpha y) - \phi_k(\lambda x)| \leq |\alpha||\phi_k(y) - \phi_k(x)| + |\alpha - \lambda||\phi_k(x)| < (|\lambda| + \delta)\delta + \delta|\phi_k(x)|.
    \]
    Считая $\delta < 1$ (без ограничения общности) получаем
    \[
        |\phi_k(\alpha y) - \phi_k(\lambda x)| \leq \delta(|\lambda| + 1 + |\phi_k(x)|) \leq \delta(|\lambda| + 1 + \max_{k}|\phi_k(x)|).
    \]
    И тогда нам достаточно взять $\delta$ меньше чем $\min\left\{1, \frac{\epsilon}{|\lambda| + 1 + \max_{k}|\phi_k(x)|}\right\}$.
\end{proof}
\chapter{Метризуемость и геометрические свойства слабых топологий.}
\section{Геометрия окрестности точки.}
\begin{note}
    Пусть $(X, \tau_{\Phi})$ -- пространство с $\Phi$-слабой топологией.
    Заметим, что
    \[
        \beta_0 = \left\{\bigcap_{k = 1}^{N} V(0, \phi_k, \epsilon) \mid \forall k \ \phi_k \in \Phi, \ \epsilon > 0\right\}
    \]
    является локальной база $\tau_{\Phi}$-окрестности нуля, так как при всяком $x \in X$
    множество
    \[
        x + \bigcap_{k = 1}^{N} V(0, \phi_k, \epsilon) = \bigcap_{k = 1}^{N} V(x, \phi_k, \epsilon),
    \]
    а значит
    \[
        \{x + \beta_0\mid x \in X\} = \beta_{\Phi}.
    \]
\end{note} % Чё?
\begin{proposition}\label{prop:infdim-subspace}
    Пусть $X$ -- бесконечномерное линейное пространство и $\tau_{\Phi}$ -- $\Phi$-слабая топология на нём.
    Тогда всякая $\Phi$-слабо открытое множество содержит бесконечномерное аффинное подпространство.
\end{proposition}
\begin{proof}
    Пусть $\phi \in \Phi$.
    Тогда, либо $\phi \equiv 0$ и тогда $\ker \phi = X$, или же $\phi \not\equiv 0$ и $\exists x_0 \in X\colon \phi(x_0) \neq 0$, а значит
    \[
        X = \ker \phi \oplus \Lin\{x_0\}.
    \]
    То есть $\forall \phi \in \Phi \ \exists x_0 \in X\colon X = \Lin\{x_0\} \oplus \ker \phi$.
    Заметим, что если мы теперь возьмём второй функционал $\phi_1$ и сузим его до $\ker \phi$, то мы опять сможем его разложить как
    \[
        \ker \phi = \Lin\{x_1\} \oplus (\ker \phi \cap \ker \phi_1),
    \]
    то есть
    \[
        X = \Lin\{x_0, x_1\} \oplus (\ker \phi \cap \ker \phi_1).
    \]
    Продолжая процедуру индуктивно, мы получим что для $\phi_1, \ldots, \phi_n$ справедливо следующее представление $X\colon$
    \[
        X = \Lin\{x_1, \ldots, x_n\} \oplus \bigcap_{k = 1}^{n} \ker \phi_k,
    \]
    при некоторых $x_1, \ldots, x_n \in X$.

    Заметим, что каждое из ядер функционалов имеет конечную коразмерность, а значит их конечное пересечение также имеет конечную коразмерность -- и, соответственно, бесконечную размерность.
    Поскольку всякое открытое множество $U$ содержит некоторое множество из базы вида $\bigcap_{k = 1}^{N} V(x, \phi_k, \epsilon)$, то есть его точки описываются как
    \[
        \forall k \hookrightarrow |\phi_k(y) - \phi_k(x)| < \epsilon,
    \]
    то в этом пересечении лежит и $x + \bigcap_{k = 1}^{N} \ker \phi_k$ -- сдвинутое бесконечномерное подпространство.
\end{proof}
\begin{corollary}\label{cor:weak-weak-star-comparison}
    Пусть $X$ -- нормированное пространство.
    \begin{enumerate}
        \item Если $\dim X = n$, то $\dim X^* = n$ и $\tau_w = \tau_n$ -- нормированная топология в $X$, а $\tau_{w*} = \tau_{n*}$ -- нормированная топология в $X^*$.
        \item Если $\dim X = \infty$, то $\tau_w \subsetneq \tau_n$.
        \item Если $\dim X = \infty$, то $\tau_{w*} \subset (\tau_w)^* \subsetneq (\tau_{n})^*$, где $(\tau_w)^*$ -- слабая топология в $X^*$, порождённая $X^{**}$.
        \item Если $\Phi(X) \neq X^{**}$, то $\tau_{w*} \subsetneq (\tau_w)^*$ (где $\Phi\colon X \to X^{**}, \ x \mapsto \Phi_x$, действующее как $\Phi_x f = f(x)$ -- каноническая изометрия Банаха).
    \end{enumerate}
\end{corollary}
\begin{proof} \leavevmode
    \begin{enumerate}
        \item Пусть $\dim X = n$ и $\{e_i\}_{i = 1}^{n}$ -- базис в $X$.
        Рассмотрим $f_k\colon X \to \Cm$ такой, что $f_k(e_m) = \delta_{km}$ (то есть координатные функционалы).
        Тогда $\forall x \in X\colon$
        \[
            x = \sum_{k = 1}^{n} x_k e_k = \sum_{k = 1}^{n} f_k(x)e_k.
        \]
        Рассмотрим произвольный $f\colon X \to \Cm$.
        Для $x \in X\colon$
        \[
            f(x) = \sum_{i = 1}^{n} x_i f(e_i) = \sum_{i = 1}^{n} f_i(x)f(e_i) = \left(\sum_{i = 1}^{n} f(e_i)f_i\right)(x),
        \]
        а значит
        \[
            f = \sum_{i = 1}^{n} f(e_i)f_i,
        \]
        и $\dim X^* = n$.
        Известно, что $\tau_w \subset \tau_n$.
        Рассмотрим $B_1(0) \subset X$ -- открытый по норме единичный шар.
        Если мы покажем, что существует такая слабая окрестность нуля $U(0)$, что $U(0) \subset B_1(0)$, то мы, очевидно, покажем обратное вложение топологий (так как вокруг любой точки открытого шара есть шар некоторого радиуса, его можно отнормировать и сдвинуть в 0, и получить искомое вложение слабой окрестности).
        Пусть
        \[
            \|x\|_e = \max\limits_{k} |x_k|.
        \]
        Тогда $\exists C > 0$ такое, что
        \[
            \forall x \in X \hookrightarrow \|x\| \leq C\|x\|_e.
        \]
        И верно следующее вложение:
        \[
            U = \left\{x \in X \mid \|x\|_e < \frac{1}{C}\right\} \subset B_1(0).
        \]
        Но множество $U$ в точности совпадает с
        \[
            U = \bigcap_{k = 1}^{n} V\left(0, f_k, \frac{1}{C}\right),
        \]
        а значит открыто в слабой топологии -- что и завершает доказательство этого пункта.

        Абсолютно аналогично находится окрестность в $X^*$ для *-слабой топологии.
        \item Пусть $\dim X = \infty$.
        Тогда и $\dim X^* = \infty$, так как $\forall x_1, \ldots, x_n \in X$ (линейно независимых) мы можем определить $L_n = \Lin\{x_1, \ldots, x_n\}$ и
        \[
            f_n \colon L_n \to \Cm, \ f_k(x_m) = \delta_{km},
        \]
        а на $L_n$ норма $\|\cdot\|$ из $X$ эквивалентна координатной норме:
        \[
            L_n \ni x = \sum_{k = 1}^{n} \alpha_k x_k,
        \]
        и
        \[
            \|x\|_e = \max\limits_{k} |\alpha_k|.
        \]
        Функционалы $|f_k(x)| \leq \|x\|_e \leq C\|x\|$ -- а значит непрерывны.
        Продолжим по Хану-Банаху и получим непрерывный на $X$ функционалы $g_1, \ldots, g_n$.
        Так как при сужении на $L_n$ они линейно независимы, то они линейно независимы в принципе и мы умеем находить набор линейно независимых векторов произвольной конечной мощности, а значит $\dim X^* = \infty$.

        Теперь, в силу утверждения о бесконечномерном подпространстве $\forall f_1, \ldots, f_n \in X^* \ \exists x_1, \ldots, x_n \in X$
        \[
            \bigcap_{k = 1}^{n} \ker f_k \oplus \Lin\{x_1, \ldots, x_n\} = X.
        \]
        И $\forall x_1, \ldots, x_n \in X \ \exists f_1, \ldots, f_n \in X^*$ такие, что
        \[
            \bigcap_{k = 1}^{n} \ker \Phi_{x_k} \oplus \Lin\{f_1, \ldots, f_n\} = X^*.
        \]
        А значит любая слабая окрестность нуля содержит такое бесконечномерное подпространство и, следовательно, неограничена по норме.
        Аналогично, любая *-слабая окрестность нуля в $X^*$ также не является ограниченной по сопряжённой норме, а значит вложения $\tau_w \subset \tau_n$ и $(\tau_{w})^* \subset (\tau_{n})^*$ являются строгими (так как ни открытый по норме шар в $X$, ни открытый по сопряжённой норме шар в $X^*$ не могут быть слабо открытыми -- они ограничены).

        Вложение $\tau_{w*} \subset (\tau_{w})^*$ очевидно, так как $\tau_{w*}$ -- слабая топология, порождённая $\Phi(X)$, а $(\tau_w)^*$ порождается всем $X^{**}$, очевидно содержащим $\Phi(X)$.
        Равенство достигается в случае если $\Phi(X) = X^{**}$, то есть если $X$ -- рефлексивно, то $\tau_{w^*} = (\tau_w)^*$.
        \item Покажем сначала техническую лемму.
        \begin{lemma}[О ядрах]
            Пусть $X$ -- линейное пространство и $f_1, \ldots, f_n\colon X \to \Cm$ -- линейные функционалы, такие, что
            \[
                \bigcap_{k = 1}^{n} \ker f_k \subset \ker f.
            \]
            Тогда $f$ линейно зависит от $f_1, \ldots, f_n$, то есть существуют $\alpha_1, \ldots, \alpha_n$ такие, что
            \[
                f = \sum_{k = 1}^{n} \alpha_k f_k.
            \]
        \end{lemma}
        \begin{proof}
            Определим линейное пространство $L$ как
            \[
                L = \left\{\begin{pmatrix*} f_1(x) \\ \ldots \\ f_n(x)
                \end{pmatrix*} \in \Cm^n \mid x \in X\right\}.
            \]
            Очевидно, что если при некоторых $x, y$ вектора $\begin{pmatrix*} f_1(x) \\ \ldots \\ f_n(x)
            \end{pmatrix*}$ и $\begin{pmatrix*} f_1(y) \\ \ldots \\ f_n(y)
            \end{pmatrix*}$ совпадают, то $x - y \in \bigcap_{k = 1}^{n} \ker f_k \subset \ker f$, то есть и $f(x) = f(y)$.
            Рассмотрим $\Psi\colon L \to \Cm$, определяемый как
            \[
                \Psi(f_1(x), \ldots, f_n(x)) = f(x).
            \]
            Определение корректно по предыдущему <<очевидно>>, и $\Psi$ очевидно линеен.
            Рассмотрим некоторый $\begin{pmatrix*} \xi_1 \\ \ldots \\ \xi_n
            \end{pmatrix*} \in L$.
            Существуют некоторые $a_1, \ldots, a_n \in \Cm$ (не зависящие от $\xi_k$) такие, что
            \[
                \Psi(\xi_1, \ldots, \xi_n) = \sum_{k = 1}^{n} a_k \xi_k.
            \]
            А значит
            \[
                \Psi(f_1(x), \ldots, f_n(x)) = \sum_{k = 1}^{n} a_k f_k(x) = f(x),
            \]
            и это верно для всякой $x \in X$, то есть
            \[
                f = \sum_{k = 1}^{n} a_k f_k,
            \]
            что и требовалось показать.
        \end{proof}
        А теперь докажем ещё одну лемму.
        \begin{lemma}[О сопряжённом пространстве]\label{lemma:about-continious-funcitonals}
            Пусть $X$ -- линейное пространство с $\Phi$-слабой топологией $\tau_{\Phi}$.
            Тогда $\Lin\{\Phi\}$ являются в точностью всеми $\tau_{\Phi}$-непрерывными функционалами.
        \end{lemma}
        \begin{proof}
            Очевидно, что $\forall \phi \in \Lin\{\Phi\}$, то есть всякий функционал, имеющий вид
            \[
                \phi = \sum_{k = 1}^{n} a_k \phi_k
            \]
            является $\tau_{\Phi}$-непрерывным (это верно просто по определению $\tau_{\Phi}$ и непрерывности суммы и умножения на скаляр в $\Cm$).

            Пусть $\phi\colon X \to \Cm$ является $\tau_{\Phi}$-непрерывным.
            Так как $\phi(0) = 0$, то
            \[
                \exists U(0) \in \tau_{\Phi}\colon \forall x \in U(0) \hookrightarrow |\phi(x)| < 1.
            \]
            По определению локальной базы $\exists \phi_1, \ldots, \phi_n \in \Phi$ и $\epsilon > 0$ такие, что
            \[
                \bigcap_{k = 1}^{n} V(0, \phi_k, \epsilon) \subset U(0).
            \]
            А значит верно следующее включение:
            \[
                U(0) \supset \bigcap_{k = 1}^{n} V(0, \phi_k, \epsilon) \supset \bigcap_{k = 1}^{n} \ker \phi_k.
            \]
            Для всякого $x \in \bigcap_{k = 1}^{n} \ker \phi_k$ и $\forall r > 0$ элемент $rx$ также лежит в пересечении ядер.
            Следовательно
            \[
                |\phi(rx)| < 1 \Rightarrow |\phi(x)| < \frac{1}{r} \rightarrow 0, r \rightarrow \infty,
            \]
            и, следовательно, $\phi(x) = 0$.
            Поэтому $\ker \phi \supset \bigcap_{k = 1}^{n} \ker \phi_k$, и, в силу леммы о ядрах, является линейной комбинацией $\phi_k$ -- а значит лежит в линейной оболочке $\Phi$.
        \end{proof}
        \begin{definition}
            Будем называть топологическим сопряжённым $X^*$ к $(X, \tau)$ пространство всех $\tau$-непрерывных линейных функицоналов.
        \end{definition}
        По лемме о сопряжённом пространстве $(X^*, \tau_{w*})^* = \Phi(X)$.
        Следовательно, если $\Phi(X) \subsetneq X^{**}$, то существует функционал из $X^{**}$, который разрывен относительно $\tau_{w*}$.
        Тогда существует открытое в $\Cm$ множество, прообраз которого открыт по $(\tau_w)^*$, но не является открытым по $\tau_{w*}$, и, следовательно, *-слабая топология на $X^*$ строго слабее слабой топологии на $X^*$.
    \end{enumerate}

\end{proof}
\begin{note}
    Если $(X, \|\cdot\|)$ -- бесконечномерное нормированное пространство, и
    \[
        S = \left\{x \in X \mid \|x\| = 1\right\},
    \]
    то слабое замыкание $S$ совпадает с замкнутым единичным шаром.

    Это действительно очевидно, так как $\forall x_0 \in X \setminus B_1[0]$ существует $f_0 \in X^*$ такой, что $\|f_0\| = 1$ и $f_0(x_0) = \|x_0\|$.
    И тогда $x_0 \in V(x_0, f_0, \|x_0\| - 1)$, а $V(x_0, f_0, \|x_0\| - 1) \cap B_1[0] = \emptyset$:
    \[
        \forall x \in V(x_0, f_0, \|x_0\| - 1) \hookrightarrow \|x\| \geq |f_0(x)| \geq |f_0(x_0)| - |f_0(x_0) - f_0(x)| > 1.
    \]
    А значит $X \setminus B_1[0] \in \tau_w$.
    С другой стороны, всякая слабая окрестность точки из открытого единичного нормированного шара содержит сдвинутое бесконечномерное подпространство, а значит точку со сколь угодно большой нормой и отрезок между центром шара и этой точкой.
    В силу теоремы о промежуточном значении она содержит и точку единичной нормы, а следовательно её пересечение со сферой непусто.
    Таким образом утверждение показано.
\end{note}
\begin{note}
    В $\ell^1$ слабая сходимость и сходимость по норме совпадают (это утверждение называется теоремой Шура).
\end{note}